\documentclass[11pt,a4paper]{article}

% ---------------------------------------------------------
% PREAMBLE (stable, warning-free, Overleaf-safe)
% ---------------------------------------------------------
\usepackage[utf8]{inputenc}
\usepackage[T1]{fontenc}
\usepackage{lmodern}
\usepackage{amsmath, amssymb, amsfonts}
\usepackage{bm}
\usepackage{geometry}
\usepackage{graphicx}
\usepackage{hyperref}
\usepackage{cite}
\usepackage{authblk}
\usepackage{physics}
\usepackage{mathtools}

\geometry{margin=2.4cm}

\title{\textbf{Companion XLVII: The Bispectrum in the Relaxation-Driven Cyclic Cosmology}}
\author{Michael Lehmann \\ \texttt{mi.lehmann@gmx.de}}
\date{April 2026}

\begin{document}
\maketitle

\begin{abstract}
The bispectrum provides a powerful probe of non-Gaussianity and mode coupling in 
the large-scale structure of the Universe. In the standard cosmological model, 
non-Gaussianity arises from non-linear gravitational evolution and baryonic 
processes, while primordial non-Gaussianity is typically assumed to be small. In 
the Relaxation-Driven Cyclic Cosmology (RDCC), the scale-dependent evolution of 
the gravitational potential introduces a new source of non-Gaussianity, as the 
relaxation mechanism modifies the temporal evolution of density perturbations in a 
scale-dependent manner. This Companion develops a continuous description of the 
bispectrum in RDCC, deriving how the relaxation scale influences the shape, 
amplitude, and configuration dependence of the bispectrum. The resulting framework 
provides a distinctive observational signature of the RDCC gravitational field.
\end{abstract}
% ---------------------------------------------------------
\section{Introduction}

The bispectrum is the Fourier transform of the three-point correlation function and 
provides a direct measure of non-Gaussianity in the matter distribution. While the 
power spectrum captures the variance of density fluctuations, the bispectrum 
captures the coupling between modes of different scales. In the standard 
cosmological model, this coupling arises from non-linear gravitational evolution 
and baryonic processes.

In RDCC, the gravitational potential evolves differently. The relaxation mechanism 
suppresses the decay of small-scale modes and enhances the coherence of large-scale 
modes. This modifies the coupling between modes in a scale-dependent manner, 
producing a distinctive bispectrum signature that reflects the relaxation scale 
$k_{\mathrm{rel}}$.
% ---------------------------------------------------------
\section{The RDCC Density Field and Mode Coupling}

The density contrast evolves according to
\begin{equation}
    \delta(\mathbf{k},a)
    = D(a)\,\delta(\mathbf{k},a_{i})
      + \int d^{3}q\,
        F_{2}(\mathbf{q},\mathbf{k}-\mathbf{q},a)\,
        \delta(\mathbf{q})\,
        \delta(\mathbf{k}-\mathbf{q})
      + \cdots,
\end{equation}
where $F_{2}$ is the second-order kernel.

In RDCC, the kernel becomes scale-dependent:
\begin{equation}
    F_{2,\mathrm{RDCC}}
    = F_{2}
      \left[
        1 - \chi_{F}
        \left(\frac{k}{k_{\mathrm{rel}}}\right)^{\alpha}
      \right].
\end{equation}

This modification alters the coupling between modes and introduces a new source of 
non-Gaussianity.
% ---------------------------------------------------------
\section{The RDCC Bispectrum}

The bispectrum is defined as
\begin{equation}
    \langle
      \delta(\mathbf{k}_{1})
      \delta(\mathbf{k}_{2})
      \delta(\mathbf{k}_{3})
    \rangle
    = (2\pi)^{3}
      \delta_{D}(\mathbf{k}_{1}
                +\mathbf{k}_{2}
                +\mathbf{k}_{3})
      B(k_{1},k_{2},k_{3}).
\end{equation}

In RDCC, the bispectrum becomes
\begin{equation}
    B_{\mathrm{RDCC}}
    = 2 F_{2,\mathrm{RDCC}}(k_{1},k_{2})
      P_{\mathrm{RDCC}}(k_{1})
      P_{\mathrm{RDCC}}(k_{2})
      + \text{cyc.},
\end{equation}
where the RDCC power spectrum is
\begin{equation}
    P_{\mathrm{RDCC}}(k)
    = P(k)
      \left[
        1 - \chi_{P}
        \left(\frac{k}{k_{\mathrm{rel}}}\right)^{\alpha}
      \right].
\end{equation}

The bispectrum therefore acquires a scale-dependent modification that reflects the 
relaxation scale.
% ---------------------------------------------------------
\section{Shape Dependence and RDCC Signatures}

The bispectrum depends on the shape of the triangle formed by the three wave 
vectors. In the standard cosmological model, the bispectrum is largest for 
squeezed, equilateral, and folded configurations.

In RDCC, the suppression of small-scale modes reduces the amplitude of the 
bispectrum in equilateral configurations, while the enhanced coherence of 
large-scale modes increases the amplitude in squeezed configurations.

The RDCC bispectrum therefore exhibits a characteristic enhancement in squeezed 
triangles, providing a direct observational probe of the relaxation scale.
% ---------------------------------------------------------
\section{The Reduced Bispectrum and Observational Probes}

The reduced bispectrum is defined as
\begin{equation}
    Q(k_{1},k_{2},k_{3})
    = \frac{
        B(k_{1},k_{2},k_{3})
      }{
        P(k_{1})P(k_{2})
        + P(k_{2})P(k_{3})
        + P(k_{3})P(k_{1})
      }.
\end{equation}

In RDCC, this becomes
\begin{equation}
    Q_{\mathrm{RDCC}}
    = Q
      \left[
        1 - \chi_{Q}
        \left(\frac{k}{k_{\mathrm{rel}}}\right)^{\alpha}
      \right].
\end{equation}

This modification produces a distinctive scale dependence that can be measured in 
galaxy surveys, weak lensing surveys, and CMB lensing reconstructions.
% ---------------------------------------------------------
\section{Conclusion}

The bispectrum in the Relaxation-Driven Cyclic Cosmology reflects the 
scale-dependent evolution of the gravitational potential. The relaxation mechanism 
modifies the mode-coupling kernel, the power spectrum, and the shape dependence of 
the bispectrum, producing distinctive signatures in squeezed, equilateral, and 
folded configurations. These effects provide a direct observational probe of the 
relaxation scale and the temporal structure of the RDCC gravitational field. The 
present Companion establishes the theoretical foundation for future studies of 
non-Gaussianity, mode coupling, and the higher-order statistics of large-scale 
structure in RDCC.

% ========================
% References
% ========================

\begin{thebibliography}{10}

\bibitem{Flagship} Lehmann, M. (2026). Relaxation-Driven Cyclic Cosmology (RDCC): A Minimal CPT-Symmetric Framework with a Single Parameter. Flagship Paper. Zenodo: \url{https://zenodo.org/records/19744958} (v43.0 or later).

\bibitem{Zenodo} The complete RDCC companion series (I--LVII) is available at the same Zenodo record above.

% Thematisch relevante Companions
\bibitem{CompXVI} Companion XVI: Boltzmann Code and Modified Hierarchy.
\bibitem{CompXXXI} Companion XXXI: RDCC Halo Model and Non-Linear Structure.
\bibitem{CompXXXII} Companion XXXII--XXXIX: Galaxy--Halo Connection, Cosmic Web, Lensing, RSD, ISW, tSZ, Rotation Curves, CGM.

% Wichtige externe Referenzen
\bibitem{Planck2020} Planck Collaboration (2020), Astron. Astrophys. 641, A6.
\bibitem{DESI2024} DESI Collaboration (2024/2025), arXiv: [aktuelle $f\sigma_8$/BAO-Paper].
\bibitem{JWST2024} Maiolino et al. (2024), Nature (JWST high-$z$ galaxies/quasars).
\bibitem{Kaiser1987} Kaiser, N. (1987), Mon. Not. R. Astron. Soc. 227, 1 (RSD).
\bibitem{Bartelmann2001} Bartelmann, M. \& Schneider, P. (2001), Phys. Rep. 340, 291 (Weak Lensing Review).
\bibitem{HuOkamoto2002} Hu, W. \& Okamoto, T. (2002), Astrophys. J. 574, 566 (CMB Lensing).
\bibitem{LewisChallinor2006} Lewis, A. \& Challinor, A. (2006), Phys. Rep. 429, 1 (CMB Lensing Review).
\bibitem{NFW1997} Navarro, J. F., Frenk, C. S. \& White, S. D. M. (1997), Astrophys. J. 490, 493 (Halo Profiles).

\end{thebibliography}

\end{document}
